\section{Gravitational slope on a uniformly rotating sphere}
\label{app:rotsphere}

This appendix derives the closed-form slope field used as a verification test in
Section~\ref{sec:numval}, and shows explicitly how it discriminates between the
two possible signs of the centrifugal term.

Consider a homogeneous sphere of radius $R$ and mass $M$, rotating uniformly at
rate $\omegarot$ about the $z$ axis, and write $g = GM/R^{2}$ for the magnitude
of the surface attraction. Because the configuration is axisymmetric, take
without loss of generality a surface point at longitude zero and latitude
$\phi$,

\begin{equation}
\mathbf{r} \;=\; R\,(\cos\phi,\;0,\;\sin\phi),
\qquad
\nhat \;=\; (\cos\phi,\;0,\;\sin\phi),
\label{eq:apppoint}
\end{equation}

the outward normal of a sphere being radial. The self-gravitational attraction
of a homogeneous sphere at its surface is radial and inward,

\begin{equation}
\ngrav \;=\; -\,g\,(\cos\phi,\;0,\;\sin\phi),
\label{eq:appattract}
\end{equation}

while the component of $\mathbf{r}$ perpendicular to the spin axis is
$\mathbf{r}_{\perp} = R\cos\phi\,(1,0,0)$, so that from Eq.~\eqref{eq:acf}

\begin{equation}
\acf \;=\; \omegarot^{2} R\cos\phi\,(1,\;0,\;0).
\label{eq:appcf}
\end{equation}

Introduce the dimensionless ratio of centrifugal to gravitational acceleration
at the equator,

\begin{equation}
k \;\equiv\; \frac{\omegarot^{2}R}{g},
\label{eq:appk}
\end{equation}

For every body considered in this paper $k$ is well below unity, and $k<1$ is
assumed in what follows unless stated otherwise; the case $k>1$ is treated
explicitly below. Adding Eqs.~\eqref{eq:appattract} and~\eqref{eq:appcf} gives
the effective acceleration

\begin{equation}
\geff \;=\; g\bigl(\,(k-1)\cos\phi,\;\;0,\;\;-\sin\phi\,\bigr).
\label{eq:appgeff}
\end{equation}

The slope of Eq.~\eqref{eq:slope} is the angle between $-\geff$ and $\nhat$.
Forming the inner product,

\begin{equation}
-\geff\cdot\nhat
\;=\; g\bigl[(1-k)\cos^{2}\phi + \sin^{2}\phi\bigr]
\;=\; g\bigl(1 - k\cos^{2}\phi\bigr),
\label{eq:appdot}
\end{equation}

and the magnitude,

\begin{equation}
\lVert\geff\rVert \;=\; g\sqrt{(1-k)^{2}\cos^{2}\phi + \sin^{2}\phi},
\label{eq:appmag}
\end{equation}

so that

\begin{equation}
\cos\slope(\phi) \;=\;
\frac{1-k\cos^{2}\phi}
     {\sqrt{(1-k)^{2}\cos^{2}\phi + \sin^{2}\phi}},
\label{eq:approtsphere}
\end{equation}

which is Eq.~\eqref{eq:rotsphere}. The factor $g$ cancels, so the slope field
depends on the body only through $k$.

Two limits check the result. At the pole, $\phi = 90^{\circ}$, both numerator
and denominator reduce to unity and $\slope = 0$: the spin axis passes through
the point, the centrifugal term vanishes there, and the effective gravity is
radial. At the equator, $\phi = 0$, the expression becomes $(1-k)/\lvert
1-k\rvert$, which is $+1$ for $k<1$ and again gives $\slope = 0$, the
centrifugal term being antiparallel to the attraction and so changing its
magnitude but not its direction. For $k>1$ the same expression gives $-1$ and
hence $\slope = 180^{\circ}$: the effective acceleration points outward and
loose material is no longer bound, which is the rotational disruption limit. The
slope is therefore zero at both extremes of latitude and attains its maximum in
between, which is what makes intermediate latitudes the informative ones for
verification.

The discriminating power of the test follows from the appearance of $k$ in
Eq.~\eqref{eq:approtsphere}. Reversing the sign of $\acf$ in
Eq.~\eqref{eq:geff}---adding the centrifugal acceleration to the attraction
rather than opposing it---is equivalent to replacing $k$ by $-k$ throughout, so
that the predicted slope becomes

\begin{equation}
\cos\slope_{\text{wrong}}(\phi) \;=\;
\frac{1+k\cos^{2}\phi}
     {\sqrt{(1+k)^{2}\cos^{2}\phi + \sin^{2}\phi}}.
\label{eq:appwrong}
\end{equation}

At $k = 0.5$ the two expressions give the values in
Table~\ref{tab:approtsphere}. They differ by $10.2^{\circ}$ at latitude
$30^{\circ}$ and by $7.1^{\circ}$ at $45^{\circ}$---differences far larger than
the discretisation error of a moderately refined mesh, so the test is decisive
rather than marginal. Both expressions vanish at $\phi = 0$ and
$\phi = 90^{\circ}$, which is precisely why a test confined to a non-rotating
body, or one that examined only the equator and poles, would detect nothing.

\begin{table}[t]
\centering
\small
\caption{Gravitational slope on a uniformly rotating sphere at
$k = \omegarot^{2}R/g = 0.5$, for the correct sign of the centrifugal term
(Eq.~\ref{eq:approtsphere}) and for the reversed sign
(Eq.~\ref{eq:appwrong}), together with the values returned by the
implementation on an icosphere mesh.}
\label{tab:approtsphere}
\begin{tabular}{rrrr}
\toprule
$\phi$ & Eq.~\eqref{eq:approtsphere} & Eq.~\eqref{eq:appwrong} & This work \\
\midrule
$30^{\circ}$ & $19.11^{\circ}$ & $8.95^{\circ}$  & $18.97^{\circ}$ \\
$45^{\circ}$ & $18.43^{\circ}$ & $11.31^{\circ}$ & $18.33^{\circ}$ \\
$60^{\circ}$ & $13.90^{\circ}$ & $10.89^{\circ}$ & $13.91^{\circ}$ \\
\bottomrule
\end{tabular}
\end{table}

The residual between the computed and analytic values, at most $0.14^{\circ}$,
is consistent with the polygonal approximation of the sphere at the mesh
resolution used and decreases under refinement.
